\chapter*{پیش‌گفتار}
\addcontentsline{toc}{chapter}{پیش‌گفتار}
\markboth{پیش‌گفتار}{پیش‌گفتار}

مایلم داستانی برایتان روایت کنم.

این روایت، شرح یک دگرگونی است؛ اما نه از آن جنسی که به نگارش نخستین هسته‌ی لینوکس توسط «لینوس توروالدز» در سال ۱۹۹۱ می‌پردازد، چرا که این تاریخچه در منابع متعددی مستند شده است. همچنین، این داستان بر بنیان‌گذاری پروژه‌ی گنو (\en{GNU}) توسط «ریچارد استالمن» در سال‌های پیش از آن تمرکز ندارد؛ هرچند آن نیز نقطه‌ی عطفی حیاتی است که در اغلب متون لینوکسی مورد بحث قرار می‌گیرد. در حقیقت، این روایت به موضوعی اساسی‌تر اختصاص دارد: بازپس‌گیری کنترل بی‌قیدوشرط بر رایانه‌ی شخصی خودتان.

زمانی که در اواخر دهه‌ی ۱۹۷۰ به عنوان دانشجو کار با رایانه‌ها را آغاز کردم، یک انقلاب فناورانه در جریان بود. ابداع ریزپردازنده این امکان را فراهم ساخت که افراد عادی، فارغ از وابستگی‌های سازمانی، مالک یک رایانه شوند. برای نسل امروز، درک دنیایی که در آن رایانه‌ها صرفاً در انحصار سازمان‌های بزرگ و دولت‌ها بودند دشوار است؛ اما باید اذعان داشت که گستره‌ی کاربردپذیری و آزادی عمل در آن دوران بسیار محدود بود.

امروزه، جهان به‌کلی دگرگون شده است. رایانه‌ها به عنصری فراگیر تبدیل شده‌اند؛ از گجت‌های پوشیدنی کوچک تا مراکز داده‌ی عظیم. یک شبکه‌ی جهانی نیز این سامانه‌ها را به یکدیگر پیوند داده است. این شرایط، فرصت‌های تازه‌ای برای یادگیری، ساختن و ابداع در اختیار افراد عادی قرار می‌دهد. با این حال، طی دهه‌های اخیر پدیده‌ای موازی نیز رخ داده است: چند شرکت بزرگ، کنترل بخش عمده‌ای از رایانه‌های جهان را در اختیار گرفته‌اند و حدود فعالیت‌های مجاز و غیرمجاز کاربران را دیکته می‌کنند.

خوشبختانه، جامعه‌ای جهانی از متخصصان و علاقه‌مندان در برابر این انحصار ایستادگی کرده‌اند. آن‌ها با توسعه‌ی نرم‌افزارهای آزاد، برای حفظ حاکمیت بر ابزارهای دیجیتال خود می‌جنگند و در این مسیر، لینوکس را تکامل می‌بخشند. بسیاری هنگام صحبت درباره‌ی لینوکس واژه‌ی «آزادی» را به کار می‌برند، اما شاید درک عمیقی از معنای واقعی آن نداشته باشند. آزادی در این بستر، یعنی اینکه خود شما تعیین کنید رایانه‌تان چگونه کار کند و تنها راه دستیابی به این هدف، شناخت عمیق از کارکرد آن سامانه است. آزادی در یک رایانه‌ی «شفاف» تجلی می‌یابد؛ سامانه‌ای که امکان کاوش و درک تمامی جزئیات آن برای کاربر مهیاست.

\section*{چرایی استفاده از رابط خط فرمان (CLI)}

آیا تا به حال به فیلم‌هایی که در آن‌ها یک هکر چیره‌دست در کمتر از ۳۰ ثانیه به سامانه‌های امنیتی نفوذ می‌کند دقت کرده‌اید؟ او هرگز از ماوس استفاده نمی‌کند! این رویکرد به این دلیل است که فیلم‌سازان به‌صورتی غریزی دریافته‌اند که تنها روش واقعی برای انجام فعالیت‌های جدی با رایانه، تعامل مستقیم اندیشه با ورودی‌های متنی از طریق صفحه‌کلید است.

اکثر کاربران امروزی صرفاً با رابط کاربری گرافیکی (\en{GUI}) آشنا هستند و تحت تأثیر الگوهای تجاری، رابط خط فرمان (\en{CLI}) را منسوخ یا هولناک می‌پندارند. این نگاه تأسف‌آور است؛ چرا که \en{CLI} ابزاری شگفت‌انگیز برای تعامل با رایانه است، درست همان‌طور که زبان نوشتاری بستر اصلی ارتباطات انسانی محسوب می‌شود. جمله‌ای ماندگار در دنیای رایانش می‌گوید: «رابط‌های گرافیکی انجام کارهای ساده را تسهیل می‌کنند، در حالی که رابط‌های خط فرمان، اجرای امور دشوار را میسر می‌سازند.»

لینوکس به دلیل ریشه در خانواده‌ی سیستم‌عامل‌های یونیکس (\en{Unix})، میراثی غنی از ابزارهای خط فرمان را به ارث برده است. یونیکس پیش از آنکه رابط‌های گرافیکی به مقبولیت عام برسند، به جایگاه برجسته‌ای دست یافت و به همین دلیل، یک رابط خط فرمان بسیار بالغ و جامع برای آن توسعه یافت. در واقع، یکی از دلایل اصلی که متخصصان اولیه، لینوکس را به جای سامانه‌هایی مانند \en{Windows NT} برگزیدند، همین \en{CLI} توانمند بود که «انجام کارهای دشوار را ممکن می‌ساخت.»

\section*{محورهای محتوایی کتاب}

این کتاب جستاری جامع و عمیق در زیست‌بوم «رابط خط فرمان» (\en{CLI}) لینوکس است. برخلاف بسیاری از متون که صرفاً بر یک ابزار خاص، نظیر پوسته \cmd{Bash}، تمرکز می‌کنند، این اثر می‌کوشد دیدگاهی کل‌نگر نسبت به فلسفه و شیوه‌ی تعامل با خط فرمان ارائه دهد. در این مسیر، نحوه‌ی عملکرد، قابلیت‌های کاربردی و بهترین شیوه‌های بهره‌گیری از این محیط به‌دقت بررسی می‌شوند.

این اثر صرفاً یک راهنمای مدیریت سیستم (\en{System Administration}) نیست. اگرچه هر بحث جدی در قلمرو خط فرمان ناگزیر به مباحث مدیریتی گره می‌خورد، اما این کتاب تنها بر مفاهیم زیرساختی و ضروری این حوزه تمرکز دارد. با این حال، با پی‌ریزی یک بنیان مستحکم در کاربری خط فرمان — که ابزاری اجتناب‌ناپذیر برای هر متخصص سیستمی است — خواننده را برای پیمودن مسیرهای تخصصی‌تر و یادگیری‌های عمیق آماده می‌سازد.

تمرکز این کتاب به‌طور مطلق بر لینوکس معطوف است. بسیاری از آثار مشابه برای جذب مخاطب گسترده‌تر، پلتفرم‌های دیگر نظیر یونیکس یا \en{macOS} را نیز پوشش می‌دهند؛ رویکردی که اغلب به قیمت کاهش عمق محتوا و بسنده کردن به مباحث سطحی تمام می‌شود. در مقابل، این کتاب منحصراً به توزیع‌های مدرن لینوکس می‌پردازد. هرچند بخش عمده‌ای از مطالب برای کاربران سایر سیستم‌عامل‌های شبه‌یونیکس نیز ثمر‌بخش است، اما هسته‌ی اصلی محتوا بر نیازهای کاربر امروزی لینوکس استوار است.

\section*{مخاطبان کتاب}

این اثر برای کاربرانی تدوین شده است که در آستانه‌ی ورود به دنیای لینوکس قرار دارند یا از پلتفرم‌های دیگر به این محیط مهاجرت کرده‌اند. مخاطب هدف ما احتمالاً یک «کاربر حرفه‌ای» (\en{Power User}) در اکوسیستم ویندوز بوده است؛ شاید مدیری هستید که وظیفه‌ی نگهداری از یک سرور لینوکسی به شما محول شده، یا شیفته‌ی دنیای رایانه‌های تک‌بردی (\en{SBC}) نظیر رزبری‌پای (\en{Raspberry Pi}) گشته‌اید. حتی ممکن است کاربر دسکتاپی باشید که در جست‌وجوی امنیت و پایداری بالاتر به لینوکس روی آورده است. در هر صورت، این کتاب راهنمای شماست.

با این وجود، باید صراحتاً بیان کرد که برای درک عمیق لینوکس هیچ میان‌بر ساده‌ای وجود ندارد. تسلط بر رابط خط فرمان فرآیندی چالش‌برانگیز و مستلزم استمرار است. این چالش نه از پیچیدگی ذاتی، بلکه از وسعت شگرف این محیط ناشی می‌شود؛ یک سیستم لینوکس استاندارد میزبان هزاران ابزار کاربردی است که از طریق \en{CLI} مدیریت می‌شوند. بنابراین، یادگیری خط فرمان را نباید یک فعالیت تفننی پنداشت.

در نقطه‌ی مقابل، دستاورد این تلاش بسیار ارزشمند است. اگر تصور می‌کنید هم‌اکنون یک کاربر توانمند هستید، باید بگوییم که هنوز اوج بهره‌وری و کاربری پیشرفته‌ی رایانه را تجربه نکرده‌اید. افزون بر این، برخلاف بسیاری از مهارت‌های نرم‌افزاری که به‌سرعت منسوخ می‌شوند، دانش خط فرمان پایدار است. مهارت‌هایی که امروز می‌آموزید، تا دهه‌های آینده نیز اعتبار و کارایی خود را حفظ خواهند کرد؛ چرا که \en{CLI} امتحان خود را در طول زمان به‌خوبی پس داده است.

در این کتاب، پیش‌فرض بر این است که شما هیچ سابقه‌ای در برنامه‌نویسی ندارید. جای نگرانی نیست؛ ما در این مسیر گام‌به‌گام همراه شما خواهیم بود و آموزش را از بنیادین‌ترین مفاهیم آغاز خواهیم کرد.

\section*{ساختار و سرفصل‌های آموزشی}

مطالب این کتاب با توالی دقیق و گزینش‌شده‌ای تدوین شده‌اند تا تجربه‌ای مشابه حضور یک مدرس خصوصی در کنار شما را تداعی کنند. بسیاری از نویسندگان، مطالب خود را به‌صورت «دانشنامه‌ای» و با پوشش طولی هر مبحث ارائه می‌دهند؛ رویکردی که اگرچه برای نگارش منطقی به نظر می‌رسد، اما برای کاربران تازه‌کار بسیار سردرگم‌کننده است. در مقابل، این اثر با رویکردی مرحله‌بندی شده، مفاهیم را متناسب با نیاز کاربر گسترش می‌دهد.

یکی دیگر از اهداف بنیادین این کتاب، آشنایی شما با «طرز تفکر یونیکسی» (\en{The Unix Way}) است که با الگوهای ذهنی رایج در محیط ویندوز تفاوت‌های اساسی دارد. در طول این مسیر، گاهی از مباحث فنی فاصله می‌گیریم تا تبیین کنیم که چرا فرآیندها به شیوه‌ای خاص عمل می‌کنند و ریشه‌های تاریخی آن‌ها چیست. لینوکس صرفاً یک سیستم‌عامل نیست، بلکه جزیی از فرهنگ دیرپای یونیکس است که زبان و هویت مختص به خود را دارد. در این میان، برای حفظ دقت علمی، حتی به برخی کاستی‌های موجود نیز اشاره خواهد شد.

این کتاب در چهار بخش اصلی انتظام یافته است که هر یک جنبه‌ای کلیدی از تجربه‌ی کار در رابط خط فرمان را پوشش می‌دهد:

\begin{enumerate}
  \item \textbf{بخش اول: یادگیری پوسته:} این بخش، نقطه آغازین کاوش ما در زبان پایه‌ی خط فرمان است. در اینجا با مفاهیم زیربنایی همچون ساختار نحو (\en{Syntax}) دستورات، پیمایش در سلسله‌مراتب سیستم فایل، فنون ویرایش در خط فرمان و روش‌های بازیابی مستندات و راهنمای ابزارها آشنا می‌شویم.
  \item \textbf{بخش دوم: پیکربندی و محیط:} در این مرحله، بر ویرایش فایل‌های پیکربندی تمرکز می‌کنیم؛ متونی که از طریق خط فرمان، رفتار و عملیات‌های سطح سیستم را تحت کنترل قرار می‌دهند.
  \item \textbf{بخش سوم: وظایف رایج و ابزارهای ضروری:} در این بخش، مجموعه‌ای از فعالیت‌های روزمره که با خط فرمان بهینه‌سازی می‌شوند، بررسی می‌گردند. سیستم‌عامل‌های شبه‌یونیکس نظیر لینوکس، میزبان برنامه‌های کلاسیک متعددی هستند که برای پردازش‌های سنگین و هوشمند داده‌ها طراحی شده‌اند.
  \item \textbf{بخش چهارم: نگارش اسکریپت‌های پوسته:} این بخش به معرفی برنامه‌نویسی پوسته می‌پردازد؛ تکنیکی بنیادین و در عین حال سهل‌الوصول برای خودکارسازی (\en{Automation}) وظایف تکراری. با یادگیری اسکریپت‌نویسی، علاوه بر افزایش بهره‌وری، با منطق برنامه‌نویسی آشنا می‌شوید که در بسیاری از زبان‌های دیگر نیز کاربرد دارد.
\end{enumerate}

\section*{روش مطالعه کتاب}

این کتاب با ساختاری خطی تدوین شده است؛ لذا پیشنهاد می‌شود مباحث را به‌صورت ترتیبی، از ابتدا تا انتها دنبال کنید. این اثر بیش از آنکه یک کتاب مرجع (\en{Reference}) ساده باشد، به شکل یک روایت آموزشی با آغاز، میانه و پایان طراحی شده است.

\section*{پیش‌نیازهای عملیاتی}

تنها پیش‌نیاز لازم برای بهره‌گیری از مفاد این کتاب، دسترسی به یک سیستم‌عامل لینوکس در حال اجرا است. شما می‌توانید این محیط را به یکی از دو روش زیر مهیا کنید:

\begin{enumerate}
  \item \textbf{نصب لینوکس روی رایانه:} انتخاب توزیع (\en{Distribution}) لینوکس کاملاً به سلیقه شما بستگی دارد، هرچند اغلب کاربران با توزیع‌های \en{Ubuntu}، \en{Fedora} یا \en{OpenSUSE} آغاز می‌کنند. اگر در انتخاب خود مردد هستید، پیشنهاد ما شروع با \en{Ubuntu} است. فرآیند نصب یک توزیع مدرن بسته به مشخصات سخت‌افزاری شما می‌تواند بسیار ساده یا تا حدی چالش‌برانگیز باشد. توصیه می‌شود از یک رایانه رومیزی با حداقل ۲ گیگابایت حافظه رم و ۶ گیگابایت فضای ذخیره‌سازی آزاد استفاده کنید. در مراحل اولیه، حتی‌الامکان از لپ‌تاپ‌ها و ارتباطات بی‌سیم (\en{Wi-Fi}) دوری کنید؛ چرا که پیکربندی درایورهای مربوط به آن‌ها معمولاً با پیچیدگی‌های بیشتری همراه است.
  \item \textbf{استفاده از سیستم زنده (\en{Live System}):} یکی از قابلیت‌های برجسته بسیاری از توزیع‌های لینوکس، امکان اجرای مستقیم آن‌ها از طریق حافظه‌های فلش یا \en{CD-ROM} بدون نیاز به نصب روی هارددیسک است. کافی است با ورود به تنظیمات \en{BIOS/UEFI}، اولویت بوت سیستم را بر روی درایو مربوطه قرار دهید. این روش راهکاری عالی برای سنجش سازگاری سخت‌افزار پیش از نصب دائمی است، هرچند ممکن است سرعت اجرای آن نسبت به نصب مستقیم روی هارددیسک کمتر باشد. توزیع‌های \en{Ubuntu} و \en{Fedora} (در کنار بسیاری از توزیع‌های دیگر) نسخه‌های زنده پایداری را ارائه می‌دهند.
\end{enumerate}

صرف‌نظر از روش انتخابی برای استقرار لینوکس، جهت اجرای برخی تمرین‌های این کتاب، گهگاه به دسترسی در سطح ابرکاربر (\en{Superuser}) یا همان امتیازات مدیریتی نیاز خواهید داشت. پس از راه‌اندازی محیط لینوکس، می‌توانید مطالعه را آغاز کرده و دستورات را به‌صورت هم‌زمان در ترمینال خود آزمایش کنید. بخش اعظمی از مطالب این کتاب ماهیت عملی دارند؛ پس پشت سیستم بنشینید و از تایپ کردن لذت ببرید!

\section*{رویکرد نگارشی درباره واژه «GNU/Linux»}

در برخی محافل تخصصی، استفاده از عنوان \en{GNU/Linux} برای اشاره به این سیستم‌عامل، از نظر ملاحظات سیاسی صحیح‌تر تلقی می‌شود. دشواری در نام‌گذاری دقیق این پلتفرم از ماهیت توسعه‌ی توزیع‌شده و مشارکت گسترده‌ی افراد بی‌شمار نشأت می‌گیرد. از منظر فنی، «لینوکس» صرفاً نام هسته (\en{Kernel}) سیستم‌عامل است، نه کل مجموعه؛ اگرچه هسته برای حیات سامانه حیاتی است، اما به تنهایی یک سیستم‌عامل کامل را تشکیل نمی‌دهد.

در این میان، نقش «ریچارد استالمن»، فیلسوف برجسته و بنیان‌گذار جنبش نرم‌افزار آزاد، بنیاد نرم‌افزار آزاد (\en{FSF}) و پروژه‌ی \en{GNU}، بسیار حائز اهمیت است. او که خالق نخستین نسخه‌ی کامپایلر \en{GNU C} و مجوز عمومی همگانی (\en{GPL}) است، اصرار دارد که به پاس خدمات پروژه‌ی \en{GNU}، این مجموعه \en{GNU/Linux} نامیده شود. با وجود پیشگامی و سهم غیرقابل‌انکار پروژه‌ی \en{GNU}، انحصار نام‌گذاری به آن می‌تواند نسبت به سایر مشارکت‌های کلیدی ناعادلانه باشد. افزون بر این، از نگاه سلسله‌مراتب سیستمی، اصطلاح \en{Linux/GNU} دقیق‌تر است؛ چرا که ابتدا هسته بارگذاری شده و سپس سایر اجزا بر روی آن اجرا می‌شوند.

در ادبیات رایج، «لینوکس» به هسته و تمامی نرم‌افزارهای آزاد و متن‌بازی اشاره دارد که یک توزیع استاندارد را شکل می‌دهند؛ یعنی کل اکوسیستم، و نه صرفاً ابزارهای \en{GNU}. از سوی دیگر، بازار سیستم‌عامل‌ها نام‌های تک‌واژه‌ای نظیر \en{DOS}، \en{Windows}، \en{macOS}، \en{Solaris} و \en{AIX} را ترجیح می‌دهد. بر همین اساس، من نیز از قالب متداول استفاده می‌کنم. با این حال، اگر شما استفاده از نام \en{GNU/Linux} را ترجیح می‌دهید، می‌توانید هنگام مطالعه، آن را در ذهن خود جایگزین کنید؛ این موضوع خدشه‌ای به محتوا وارد نخواهد کرد.

\section*{منابع مطالعه تکمیلی}

درباره شخصیت‌های کلیدی مطرح‌شده در بخش‌های پیشین:

\textbf{لینوس تروالدز (\en{Linus Torvalds}):}
\begin{latin}
\url{http://en.wikipedia.org/wiki/Linux_Torvalds}
\end{latin}

\textbf{ریچارد استالمن (\en{Richard Stallman}):}
\begin{latin}
\url{http://en.wikipedia.org/wiki/Richard_Stallman}
\end{latin}

درباره بنیاد نرم‌افزار آزاد (\en{Free Software Foundation}) و پروژه گنو (\en{GNU Project}):
\begin{latin}
\url{http://en.wikipedia.org/wiki/Free_Software_Foundation}\\
\url{http://www.fsf.org}\\
\url{http://www.gnu.org/}
\end{latin}

دیدگاه‌های استالمن پیرامون مناقشه نام‌گذاری \en{GNU/Linux}:
\begin{latin}
\url{http://www.gnu.org/gnu/why-gnu-linux.html}\\
\url{http://www.gnu.org/gnu/gnu-linux-faq.html#tools}
\end{latin}